\documentclass{resume} % Use the custom resume.cls style

\usepackage{marvosym}
\usepackage[left=0.4 in,top=0.4in,right=0.4 in,bottom=0.4in]{geometry} % Document margins

\usepackage{graphicx}
\usepackage{xcolor}
% \usepackage{emoji}
% \usepackage{emoji}

\newcommand{\tab}[1]{\hspace{.2667\textwidth}\rlap{#1}} 
\newcommand{\itab}[1]{\hspace{0em}\rlap{#1}}
\name{Mariano Maisonnave} % Your name
% You can merge both of these into a single line, if you do not have a website.
\address{+1(902) 448-5124 \\ Nova Scotia, Canada}
\address{
\href{https://www.linkedin.com/in/mariano-maisonnave/}{LinkedIn/in/mariano-maisonnave/} \\
\href{mailto:maisonnavemariano@gmail.com}{maisonnavemariano@gmail.com} \\ 
\href{https://cs.dal.ca/~maisonnave/google-scholar}{Google Scholar} 
} 


% \\ \href{https://mmaisonnave.github.io/}{mmaisonnave.github.io}
% }     %
% \address{\color{gray}(he/him)}

\begin{document}

%----------------------------------------------------------------------------------------
%	OBJECTIVE
%----------------------------------------------------------------------------------------
\vspace{-0.3cm}
\begin{rSection}{}
\textbf{Postdoctoral researcher} with over 8 years in Machine Learning, based in Halifax and authorized to work in Canada (PR status). Bring 6+ years of expertise in Natural Language Processing projects, 4+ years in causality and time series analysis for energy demand forecasting and 1 year in healthcare analytics.

\end{rSection}





\vspace{-0.2cm}
\begin{rSection}{SELECTED SKILLS}
% \vspace{-0.2cm}
\textbf{Programming Languages and tools:} Proficient with \texttt{Python}, \texttt{Java}, \texttt{C}, and \texttt{Git}.
\vspace{-0.25cm}

\textbf{Python Libraries:} Highly skilled with all major \texttt{Python} libraries (such as \texttt{NumPy}, \texttt{pandas}, \texttt{Matplotlib}).
\vspace{-0.25cm}

\textbf{Machine Learning:} Experienced with ML and LLMs libraries (\texttt{transformers},\texttt{spaCy},\texttt{PyTorch},\texttt{sklearn}, and others).
\vspace{-0.75cm}

\textbf{Sysadmin:} Skilled \texttt{Linux} enthusiast experienced in remote computing resources (\texttt{SSH/SFTP}, \texttt{Bash}, \texttt{AWS}, \texttt{Slurm}).
\end{rSection}

\vspace{-0.3cm}
\begin{rSection}{Education}

{\bf Ph.D. in Computer Science}, GPA 3.94/4, Universidad Nacional del Sur, Argentina \hfill {2017 --- 2021}\\
\underline{Thesis}: Variable selection and causal discovery from  newspaper article texts.\\
% Relevant courses: Econometrics, 

\vspace{-0.6cm}


{\bf B.Sc in Computer Systems Engineering}, GPA 3.51/4, Universidad Nacional del Sur, Argentina \hfill {2011 --- 2016}\\
% \underline{Final project}: Feature engineering techniques evaluated on supervised learning methods.
%Minor in Linguistics \smallskip \\
%Member of Eta Kappa Nu \\
%Member of Upsilon Pi Epsilon \\
\vspace{-0.2cm}


\end{rSection}

\vspace{-0.4cm}
\begin{rSection}{EXPERIENCE}




% CAPE BRETON
\includegraphics[height=3\fontcharht\font`\B]{cbu_logo.png}%
\raisebox{1.\height}{\textbf{Postdoctoral Fellow, Healthcare Analytics, Cape Breton University}} \hfill \textbf{Aug 2023 --- Aug 2024}\vspace{-0.2cm}

% \textbf{Postdoc Fellow  } \hfill Aug 2023 --- Aug 2024\\\underline{Cape Breton University} \hfill \textit{Sydney, NS}\\
% \vspace{-0.7cm}

\begin{itemize}
    \item Developed explainable ML models using 600K+ hospital discharge records from Nova Scotia to predict 30-day Unplanned Hospital Readmission (UHR) risk. \vspace{-0.2cm}
    \item Identified key risk factors and created guidelines to reduce screenings by 50\% while capturing 72\% of UHR cases.\vspace{-0.2cm}
    \item Utilized SHAP, Permutation Feature Importance, and Random Forest for model interpretability and accuracy.\vspace{-0.2cm}
    \item Handled dataset imbalance using various techniques and combined discharge records for the same patients to
    enhance detection performance.
\end{itemize}
% Our research focused on using explainable ML applied to the real-world data in the healthcare domain to identify risk factors and obtain insight for medical personnel to handle the unplanned hospital readmission problem in Nova Scotia. I handle the entire lifecycle of our research. From formulating research questions, conducting thorough literature reviews, executing experiments to produce both theoretically grounded and applicable solutions, and writing scientific articles and presenting our results in relevant venues.

% \begin{itemize}
%     \item Research focused on explainable ML applied to real-world healthcare data from Nova Scotia, Canada.
%     \vspace{-0.35cm}
%     \hspace{-1cm} \item  Aim: Identify risk factors and insights for medical personnel to handle unplanned hospital readmissions.
%     \vspace{-0.35cm}
%     \item Responsibilities:
%     \vspace{-0.35cm}
%     \begin{itemize}
%             \item Formulating research questions.
%             \vspace{-0.35cm}
%             \item Conducting thorough literature reviews.
%             \vspace{-0.35cm}
%             \item Executing experiments to produce theoretically grounded and applicable solutions.
%             \vspace{-0.35cm}
%             \item Writing scientific articles and presenting results in relevant venues.
%     \end{itemize}
% \end{itemize}


% \begin{itemize}
%     \item \textit{\textcolor{darkgray}{Mariano Maisonnave, Enayat Rajabi, Majid, and Peter, 2024. Alternate Level of Care Patients in Canada: A Systematic Review: Alternate Level of Care Patients in Canada. Accepted for publication in Canadian Geriatrics Journal.}}
% \end{itemize}
% \vspace{-0.2cm}
% \textbf{Postdoc Fellow  } \hfill Jan 2022 --- July 2023\\
% \underline{Dalhousie University} \hfill \textit{Halifax, NS}\\
\includegraphics[height=3\fontcharht\font`\B]{dal_logo.png}%
\raisebox{1.\height}{\textbf{ Postdoctoral Fellow, NLP Researcher, Dalhousie University}} \hfill \textbf{Jan 2022 --- Jul 2023}\vspace{-0.2cm}


\begin{itemize}
   \item Designed and implemented an interactive active learning NLP system in Python, enabling real-time training of a logistic regression model and facilitating user interaction for efficient information retrieval.\vspace{-0.2cm}
   \item Allowed navigation through a corpus of over 2.9 million news articles, enhancing research capabilities for social scientists.\vspace{-0.2cm}
   \item Employed large language models (Top2Vec) to categorize retrieval results by topic.
\end{itemize}

% I reviewed the state-of-the-art active learning for high-recall information retrieval to provide a solution to social science researchers conducting their research using large amounts of textual data. I designed and implemented an active learning system in Python that enabled researchers to interact with a ML model and train it in real-time. Social science researchers could perform information retrieval over millions of news articles using our system, which led them to conduct novel research in their fields.
% \begin{itemize}
%     \item Reviewed state-of-the-art active learning for high-recall information retrieval.
    
%     \item Aim: Provide a solution for social science researchers using large amounts of textual data.
%     \item Responsibilities:
%     \item Designed and implemented an active learning system in Python.
%     \item Enabled researchers to interact with a ML model and train it in real-time.
%     \item System allowed information retrieval over millions of news articles.
%     \item Facilitated novel research in social science fields.
% \end{itemize}


% \textbf{Research Fellow} \hfill Apr 2017 --- Oct 2021\\
% \underline{Institute for Computer Science and Engineering} \hfill \textit{Bahía Blanca, Argentina}\\
% \vspace{-0.7cm}

\includegraphics[height=2.5\fontcharht\font`\B]{icic.png}%
\raisebox{1.\height}{\textbf{ NLP Research Fellow, Institute for Computer Science and Engineering}} \hfill \textbf{Apr 2017 --- Oct 2021}\vspace{-0.2cm}
\begin{itemize}
    \item Led an interdisciplinary project that integrated NLP and causal inference to extract causal graphs from a corpus of 1.8 million New York Times news articles, utilizing causality learning and time series analysis tools.\vspace{-0.2cm}
    \item The first step involved extracting information from the text using deep learning techniques, specifically LSTM and CNN neural networks, built upon representations created with large language models (BERT).\vspace{-0.2cm}
    \item The second step transformed the extracted information into time series data and applied causal discovery tools, including PCMCI, PC, DirectLiNGAM, and Vector Autoregression (VAR).\vspace{-0.2cm}
    \item Combined the entire pipeline to summarize large volumes of text on specific topics into causal graph visualizations that depict relevant events and their causal interactions.
\end{itemize}

% I collaborated closely with Computer Science and Economics researchers to design a framework to extract causal graphs from a large corpus of news article texts. I developed a term-weighting technique to identify relevant concepts present in texts. I built and evaluated a ML model to identify ongoing events from news articles. We used causal learning tools on top of the extracted variables (concepts and events) to condense the information from the text into a single causal graph. As a case study, we also applied causal learning tools to a real-world energy demand time-series dataset to understand the role of causality in forecasting analysis. 

% I used causality learning and time series analysis tools (primarily in Python, some in R), natural language processing, and ML tools to conduct the project. 
% We published articles on term weighting, causal learning, natural language processing, and energy forecasting in several high-impact journals and conferences.





\includegraphics[height=3\fontcharht\font`\B]{uns_logo.png}%
\raisebox{1.\height}{\textbf{ University Instructor, Computer Science Courses}} \hfill \textbf{Oct 2015 --- Dec 2021}\vspace{-0.2cm}
\begin{itemize}
    \item Teaching assistant, head of teaching assistant, and professor at several university undergraduate courses.
    \vspace{-0.25cm}
    \item \textbf{Courses}: Object-oriented Programming, Computer Networking, and Embedded Systems. 
\end{itemize}
%  \begin{itemize}
%     \itemsep -3pt {} 
%      \item Work as teaching assistant, head of teaching assistant and professor of several undergrad courses.
%      \item Courses: computer networking, embedded systems and introduction to object-oriented programming. 

%  \end{itemize}
\end{rSection} 
%----------------------------------------------------------------------------------------
%	EDUCATION SECTION
%----------------------------------------------------------------------------------------




%----------------------------------------------------------------------------------------
% TECHINICAL STRENGTHS	
%----------------------------------------------------------------------------------------




% \begin{rSection}{SELECTED PUBLICATIONS}
% \vspace{-1.25em}
% \item \textbf{Maisonnave, M}., Delbianco, F., Tohmé, F., Milios, E. and Maguitman, A.G., 2022. \textit{Causal graph extraction from news: a comparative study of time-series causality learning techniques}. PeerJ Computer Science, 8, p.e1066. \href{https://peerj.com/articles/cs-1066/}{[link]}
% \vspace{-0.2cm}
% \item \textbf{Maisonnave, M}., Delbianco, F., Tohmé, F., Maguitman, A. and Milios, E., 2022. \textit{Detecting ongoing events using contextual word and sentence embeddings}. Expert Systems with Applications, 209, p.118257. \href{https://www.sciencedirect.com/science/article/pii/S0957417422013975}{[link]}
% \vspace{-0.2cm}
% \item \textbf{Maisonnave, M}., Delbianco, F., Tohmé, F. and Maguitman, A., 2021. \textit{Assessing the behavior and performance of a supervised term-weighting technique for topic-based retrieval}. Information Processing \& Management, 58(3), p.102483. \href{https://www.sciencedirect.com/science/article/pii/S0306457320309729}{[link]}
% % \vspace{-0.15cm}
% \vspace{-0.2cm}
% \end{rSection} 
%----------------------------------------------------------------------------------------
%	WORK EXPERIENCE SECTION
%----------------------------------------------------------------------------------------

% \begin{rSection}{PROJECTS}
% \vspace{-1.25em}
% \item \textbf{Hiring Search Tool.} {Built a tool to search for Hiring Managers and Recruiters by using ReactJS, NodeJS, Firebase and boolean queries. Over 25000 people have used it so far, with 5000+ queries being saved and shared, and search results even better than LinkedIn! \href{https://hiring-search.careerflow.ai/}{(Try it here)}}
% \item \textbf{Short Project Title.} {Build a project that does something and had quantified success using A, B, and C. This project's description spans two lines and also won an award.}
% \item \textbf{Short Project Title.} {Build a project that does something and had quantified success using A, B, and C. This project's description spans two lines and also won an award.}
% \end{rSection} 

%----------------------------------------------------------------------------------------
% \begin{rSection}{Extra-Curricular Activities} 
% \begin{itemize}
%     \item 	Actively write \href{https://www.faangpath.com/blog/}{blog posts} and social media posts (\href{https://www.tiktok.com/@faangpath}{TikTok}, \href{https://www.instagram.com/faangpath/?hl=en}{Instagram}) viewed by over 20K+ job seekers per week to help people with best practices to land their dream jobs. 
%     \item	Sample bullet point.
% \end{itemize}


% \end{rSection}

%----------------------------------------------------------------------------------------


% \newpage
% \begin{rSection}{PUBLICATIONS} 
% \begin{itemize}    
%     \item Maisonnave, M., Delbianco, F., Tohmé, F., Milios, E. and Maguitman, A., 2024. Learning causality structures from electricity demand data. Energy Systems, pp.1-23.
%     \vspace{-0.15cm}
%     \item Maisonnave, M., Delbianco, F., Tohme, F., Milios, E. and Maguitman, A.G., 2022. Causal graph extraction from news: a comparative study of time-series causality learning techniques. PeerJ Computer Science, 8, p.e1066.
%     \vspace{-0.15cm}
%     \item Ramirez-Orta, J.A., Sabando, M.V., Maisonnave, M. and Milios, E.E., 2022. MALNIS at IberLEF-2022 DETESTS Task: A Multi-Task Learning Approach for Low-Resource Detection of Racial Stereotypes in Spanish. In IberLEF@ SEPLN.
%     \vspace{-0.15cm}
%     \item Maisonnave, M., Delbianco, F., Tohmé, F., Maguitman, A. and Milios, E., 2022. Detecting ongoing events using contextual word and sentence embeddings. Expert Systems with Applications, 209, p.118257.
%     \vspace{-0.15cm}
%     \item Maisonnave, M., Delbianco, F., Tohmé, F. and Maguitman, A., 2021. Assessing the behavior and performance of a supervised term-weighting technique for topic-based retrieval. Information Processing \& Management, 58(3), p.102483.
%     \vspace{-0.15cm}
%     \item Maisonnave, M., Delbianco, F., Tohmé, F., Maguitman, A.G. and Milios, E.E., 2020, September. Assessing causality structures learned from digital text media. In Proceedings of the ACM Symposium on Document Engineering 2020 (pp. 1-4).
%     \vspace{-0.15cm}
%     \item Maisonnave, M., Delbianco, F., Tohmé, F.A. and Maguitman, A.G., 2019. A flexible supervised term-weighting technique and its application to variable extraction and information retrieval. Inteligencia Artificial, 22(63), pp.61-80.

%     % \item Recipient of a five-year scholarship to pursue a doctorate in computer science, 2017-2022. Granted by the Argentinian National Scientific and Technical Research Council (CONICET).
%     % \vspace{-0.2cm}     
%     % \item Recipient of the Encouragement of Scientific Vocations Scholarship (EVC-CIN), 2015 and 2016. Funding for undergrad students to conduct a part-time research project.
% \end{itemize}
% \end{rSection}






\vspace{-0.2cm}
\begin{rSection}{AWARDS AND SCHOLARSHIPS} 
\begin{itemize}
    \item Google Latin America Research Award (LARA) Winner: Awarded three consecutive years (2019, 2020, 2021).
    \vspace{-0.25cm}
    \item Emerging Leaders in the Americas Program (ELAP) Scholarship Recipient, 2018.
    % \item Recipient of a five-year scholarship to pursue a doctorate in computer science, 2017-2022. Granted by the Argentinian National Scientific and Technical Research Council (CONICET).
    % \vspace{-0.2cm}     
    % \item Recipient of the Encouragement of Scientific Vocations Scholarship (EVC-CIN), 2015 and 2016. Funding for undergrad students to conduct a part-time research project.
\end{itemize}
\end{rSection}


% %----------------------------------------------------------------------------------------
% \begin{rSection}{Leadership} 
% \begin{itemize}
%     \item Admin for the \href{https://discord.com/invite/WWbjEaZ}{FAANGPath Discord community} with over 6000+ job seekers and industry mentors. Actively involved in facilitating online events, career conversations, and more alongside other admins and a team of volunteer moderators! 
% \end{itemize}


% \end{rSection}


\end{document}
